\chapter{Cơ sở lý thuyết \& Các công trình liên quan}
\label{chap:co-so-ly-thuyet}
Trong chương này, chúng ta sẽ đi sâu vào hai nền tảng lý thuyết cốt lõi của nghiên cứu này: Thang phân loại nhận thức Bloom và Hệ thống đa tác tử. Phần cuối chương sẽ tổng hợp các công trình liên quan để chỉ ra khoảng trống nghiên cứu hiện tại mà khóa luận này hướng tới giải quyết.


\section{Phân loại Bloom}

Để các hệ thống sinh câu hỏi trắc nghiệm tự động có thể vừa đảm bảo các mục tiêu sư phạm vừa không biến thành một cỗ máy tạo câu hỏi ngẫu nhiên, việc thiết kế độ khó cần dựa trên một thang đo đã được kiểm chứng bởi khoa học. % [MODIFY]: Them in nghieng cho Bloom
Nghiên cứu này tập trung vào khung nhận thức được công nhận rộng rãi nhất trong lĩnh vực giáo dục là \textbf{phân loại Bloom}. Thang đo này được nhà tâm lý học và giáo dục \textit{Benjamin Bloom} đề xuất lần đầu vào năm 1956 \cite{bloom1956taxonomy} và được tinh chỉnh bởi \textit{Anderson} và \textit{Krathwohl} vào năm 2001 \cite{anderson2001taxonomy}. Trong nghiên cứu này, nhằm tối ưu hóa và phù hợp với hình thức thi câu hỏi trắc nghiệm đa lựa chọn (MCQ - Multiple-Choice Question), sáu cấp độ ban đầu của thang đo Bloom được phân loại lại thành ba cấp độ cốt lõi là Truy xuất, Suy luận và tổng hợp, và Lập luận phản biện.

\subsection{Cấp độ 1 - Truy xuất}

Cấp độ Truy xuất tương ứng với bậc Nhớ trong thang đo Bloom sửa đổi. Cấp độ này là nền tảng của sự nhận thức, tập trung chủ yếu vào khả năng ghi nhớ cơ học và truy xuất thông tin dài hạn, hoặc nhận dạng một đối tượng, sự kiện được trình bày rõ ràng trong một đoạn văn bản thuộc bài đọc. Quá trình này thường được gọi là nhận thức nông.

Trong quá trình tinh chỉnh các chỉ dẫn cho LLMs sinh câu hỏi ở cấp độ Truy xuất, hệ thống được cài đặt để tập trung vào các chi tiết cụ thể, định nghĩa thuật ngữ và nhận dạng các sự kiện. Đặc thù của MCQ ở dạng này thường xoay quanh các câu hỏi “Ai?”, “Cái gì?”, “Ở đâu?” và “Khi nào?”. Đối với bài toán  này, LLMs thể hiện vô cùng xuất sắc trong việc tự động tạo ra các câu hỏi, bởi vì quy trình này cơ bản chỉ đòi hỏi khả năng đối sánh chuỗi và trích xuất các thực thể từ văn bản mà không cần hiểu sâu hay xây dựng bất kỳ chuỗi suy luận logic nào. Không chỉ vậy, các phương án nhiễu ở mức độ này cũng tương đối dễ xây dựng; chúng thường là các thực thể hoặc dữ liệu xuất hiện trong văn bản gốc để tạo ra sự khó khăn cho người đọc.

\subsection{Cấp độ 2 - Suy luận và tổng hợp}

Cấp bậc thứ hai bao gồm hai kỹ năng Suy luận và Tổng hợp, đại diện cho các bậc Hiểu, Áp dụng và Sáng tạo trong phân loại Bloom. Ở cấp độ này, thay vì được hỏi các câu hỏi có thể truy xuất trực tiếp từ bề mặt của văn bản, người học sẽ buộc phải xử lý thông tin đã học và ứng dụng lý thuyết đó vào một ngữ cảnh hoàn toàn mới hoặc một tình huống giả định.

Trong bối cảnh bài toán MCQ, Suy luận yêu cầu người học phải có khả năng liên kết các mệnh đề nằm rải rác trong bài đọc để rút ra các kết quả, nhận định không được tác giả viết ra một cách trực tiếp. Đối với Tổng hợp, cấp độ này yêu cầu người học phải chắt lọc, sắp xếp lại các yếu tố riêng lẻ thành một góc nhìn bao quát hoặc dự đoán một kết quả tương lai. Bài toán này đòi hỏi cần có các chỉ dẫn cụ thể và hợp lý cho LLMs. Ngoài ra, việc tạo ra các đáp án nhiễu ở cấp độ này cũng là một thách thức kỹ thuật đối với AI: các đáp án sai phải được thiết kế dựa vào sự suy diễn sai lầm về mặt logic hoặc dựa vào việc người đọc hiểu sai một quy luật nào đó để từ đó có thể tạo ra các lựa chọn bẫy hợp lý về mặt trực giác. 

\subsection{Cấp độ 3 - Lập luận phản biện}

Cấp độ cuối cùng của nghiên cứu này là Cấp độ Lập luận phản biện, tương ứng với bậc Phân tích và Đánh giá. Trọng tâm của cấp độ nhận thức này là khả năng chia nhỏ một lượng thông tin đồ sộ, phức tạp thành các thành phần nhỏ hơn nhằm phân tích, nhận diện các giả định ngầm và đưa ra những nhận xét, quan điểm khách quan dựa trên sự đối chiếu với các bằng chứng thực tiễn.

Các câu hỏi trắc nghiệm ở cấp độ này thường yêu cầu người học tưởng tượng ra các tình huống kiểm tra ngụy biện logic, yêu cầu họ phân biệt rõ ràng giữa một sự thật khách quan và một quan điểm chủ quan, hoặc được yêu cầu đánh giá độ chính xác, hợp lý của một lập luận khoa học. Thách thức lớn nhất mà các LLM phải đối mặt ở cấp độ này là việc sinh ra các câu hỏi chất lượng cao chỉ có duy nhất một đáp án đúng. Quan trọng hơn cả, phương án nhiễu ở mức độ này không thể chỉ là những thông tin sai lệch thông thường, mà phải là các lập luận, lý lẽ nhìn bề mặt có vẻ cực kỳ thuyết phục nhưng lại chứa các lỗ hổng về mặt logic, từ đó buộc người học phải sử dụng tư duy phản biện cao mới có thể nhận diện được.

\section{Hệ thống đa tác tử}

\subsection{Khái niệm hệ thống đa tác tử}

Hệ thống đa tác tử đại diện cho một thế hệ kiến trúc trí tuệ nhân tạo hiện đại, cho phép giải quyết các bài toán phức tạp mà một LLM duy nhất khó có thể xử lý hiệu quả. Cách tiếp cận này chia nhỏ bài toán thành các nhiệm vụ hẹp hơn và giao cho một mạng lưới gồm nhiều tác tử chuyên biệt có khả năng tương tác qua lại. Một sai lầm phổ biến của nhiều hệ thống là yêu cầu một LLM giải quyết trọn vẹn một bài toán đòi hỏi quy trình phức tạp về logic và nghiệp vụ, khiến kết quả đầu ra kém chất lượng và khó sử dụng. Nguyên nhân chính là mô hình bị giao quá nhiều vai trò cùng lúc như đọc hiểu văn bản dài, sinh câu hỏi và tạo đáp án nhiễu trong cùng một lần gọi.

Hệ thống đa tác tử giải quyết những vấn đề này bằng cách mô phỏng cách con người phân công vai trò và phối hợp xử lý công việc, đồng thời dựa trên triết lý của \textbf{ReAct} \cite{yao2022react}. Thay vì chỉ trả lời một cách thụ động, khung ReAct buộc LLM thực hiện một vòng lặp nội tại gồm suy nghĩ, hành động, quan sát kết quả và điều chỉnh suy nghĩ. Trong hệ thống đa tác tử, mỗi tác tử được trang bị một chỉ thị hẹp nhưng có mức chuyên môn hóa cao, cho phép tập trung tài nguyên nhận thức vào một vấn đề duy nhất. Ví dụ, Tác tử Phân tích chỉ chịu trách nhiệm lọc và cấu trúc hóa văn bản thô; sau đó kết quả được chuyển cho các Tác tử Sinh câu hỏi để tạo ra bộ câu hỏi ban đầu, rồi tiếp tục được chuyển cho các Tác tử Đánh giá và Sửa lỗi để thực hiện vòng cải thiện chất lượng. Sự chuyên môn hóa này không chỉ giảm thiểu ảo giác mà còn giúp cải thiện chất lượng nội dung sinh ra, giảm lỗi và duy trì độ ổn định cho toàn hệ thống.

\subsection{Thư viện LangGraph}

% [MODIFY]: Them in nghieng cho LangGraph va LangChain
Để cài đặt và triển khai hệ thống đa tác tử một cách đúng đắn, nghiên cứu này sử dụng \textbf{LangGraph} \cite{langgraph2024}, một thư viện mã nguồn mở chuyên dụng được xây dựng trên nền tảng của hệ sinh thái \textbf{LangChain} \cite{langchain2023}. Khác với triết lý xử lý tuần tự tuyến tính của LangChain, LangGraph được thiết kế riêng cho việc tạo ra các luồng tác tử thông qua đồ thị có hướng kèm theo một đối tượng trạng thái dùng để lưu trữ dữ liệu khi hệ thống duyệt qua từng nút trong đồ thị.

Ba thành phần chính của một hệ thống đa tác tử khi được triển khai bằng LangGraph bao gồm:

\begin{enumerate}

    \item \textbf{Trạng thái: }Là một cấu trúc dữ liệu trung tâm, hoạt động như bộ nhớ chung để duy trì ngữ cảnh cho toàn bộ hệ thống, ví dụ như các ý chính đã phân tích từ văn bản gốc, các câu hỏi đã tạo và lỗi của từng câu hỏi.
    
    \item \textbf{Nút:} Đây là các hàm Python có nhiệm vụ xử lý logic cho từng tác tử (Nút sinh câu hỏi cho cấp độ 3, Nút thực hiện đánh giá câu hỏi cấp độ 2, …). Nút này sẽ nhận được Trạng thái hiện tại, gửi lệnh đến LLM để xử lý và sau đó cập nhật Trạng thái mới với kết quả vừa tạo ra.
    
    \item \textbf{Cạnh:} Thành phần này có nhiệm vụ xác định hướng di chuyển của luồng dữ liệu giữa các nút trong đồ thị. LangGraph cho phép hệ thống tự đưa ra quyết định rẽ nhánh, kích hoạt cơ chế thực thi song song, hoặc kết thúc luồng hoạt động dựa trên các tiêu chí xuất hiện trong thời gian thực.

\end{enumerate}

Khả năng hỗ trợ cấu trúc đồ thị của LangGraph cho phép triển khai các hệ thống trí tuệ nhân tạo tự định hướng. Nếu một tác tử Đánh giá phát hiện câu hỏi sinh ra có chứa phương án nhiễu vô lý hay nhiều đáp án đúng, hệ thống có thể yêu cầu một Tác tử Sửa lỗi xử lý vấn đề rồi đưa kết quả đó quay lại để đánh giá tiếp. Chu trình này chỉ dừng khi hệ thống tạo ra kết quả đạt chất lượng mong muốn. Ngoài ra, LangGraph còn được trang bị bộ nhớ chốt chặn giúp hệ thống có thể phục hồi trạng thái khi gặp các vấn đề như lỗi giao diện lập trình ứng dụng.

\section{Phương pháp đánh giá bằng LLM}

Để khắc phục các hạn chế về chi phí và khả năng mở rộng của đánh giá bởi con người, cộng đồng nghiên cứu đã triển khai phương pháp sử dụng mô hình ngôn ngữ lớn làm giám khảo tự động (LLM-as-a-Judge). Ý tưởng cốt lõi của phương pháp này là tận dụng năng lực hiểu ngôn ngữ và suy luận của các LLM thương mại mạnh mẽ để thay thế vai trò của con người trong các quy trình thẩm định chất lượng.

Một trong những khung đánh giá tiêu biểu nhất là \textbf{G-Eval} \cite{liu2023geval}, được đề xuất bởi \textit{Liu} và cộng sự vào năm 2023. G-Eval thiết kế một quy trình đánh giá có cấu trúc, trong đó LLM giám khảo nhận đầu vào gồm ba thành phần: nội dung nguồn, nội dung cần đánh giá và bộ tiêu chuẩn đánh giá chi tiết. Điểm cải tiến then chốt của G-Eval so với các phương pháp đánh giá tự động trước đó là yêu cầu bắt buộc tích hợp kỹ thuật chuỗi suy nghĩ (CoT - Chain-of-Thought \cite{wei2022chain}): trước khi đưa ra điểm số, LLM giám khảo phải trình bày chuỗi lập luận logic giải thích quá trình đánh giá. Cơ chế này giúp tăng tính minh bạch và cho phép con người kiểm tra tính hợp lý của kết quả đánh giá.

Phương pháp LLM-as-a-Judge mang lại ba lợi thế quan trọng. Thứ nhất, khả năng mở rộng vượt trội: hệ thống có thể đánh giá hàng nghìn câu hỏi trong thời gian ngắn với chi phí thấp hơn nhiều so với đánh giá bởi con người. Thứ hai, tính nhất quán cao: cùng một bộ tiêu chuẩn đánh giá, LLM giám khảo luôn đưa ra kết quả ổn định và có thể tái lập, loại bỏ yếu tố chủ quan giữa các người đánh giá khác nhau. Thứ ba, tính linh hoạt: bộ tiêu chuẩn đánh giá có thể được điều chỉnh dễ dàng cho các miền kiến thức và ngữ cảnh giáo dục khác nhau mà không cần đào tạo lại đội ngũ chuyên gia.

Tuy nhiên, phương pháp này cũng tồn tại các hạn chế cần lưu ý. LLM giám khảo có thể mắc các thiên lệch hệ thống như ưu tiên các câu trả lời dài hơn, thiên vị đầu ra của chính mình, hoặc thiếu nhạy cảm với các lỗi sư phạm tinh tế mà chuyên gia giáo dục có thể nhận ra ngay lập tức. Do đó, kết quả từ LLM-as-a-Judge nên được xem như một công cụ sàng lọc hiệu quả, nhưng vẫn cần được đối chiếu với đánh giá của con người để xác nhận độ tin cậy.

\section{Các công trình liên quan}

Phần này trình bày tổng quan các hướng nghiên cứu liên quan đến bài toán sinh câu hỏi trắc nghiệm tự động từ ảnh tài liệu, tập trung vào hai nhóm chính: phương pháp sinh câu hỏi truyền thống và sinh câu hỏi dựa trên mô hình ngôn ngữ lớn.

\subsection{Sinh câu hỏi trắc nghiệm tự động bằng quy tắc}

Các nghiên cứu tiên phong về sinh câu hỏi trắc nghiệm tự động chủ yếu dựa trên các kỹ thuật xử lý ngôn ngữ tự nhiên (NLP - Natural Language Processing) truyền thống và các quy tắc ngôn ngữ học được thiết kế thủ công. \textit{Heilman} và \textit{Smith} \cite{heilman2010good} đề xuất phương pháp tiếp cận dựa trên xếp hạng thống kê cho bài toán sinh câu hỏi, trong đó hệ thống sinh ra một lượng lớn câu hỏi ứng viên từ văn bản đầu vào thông qua các phép biến đổi cú pháp, sau đó sử dụng mô hình xếp hạng thống kê để lọc và chọn ra các câu hỏi có chất lượng cao nhất. Cách tiếp cận này đã trở thành một khuôn mẫu phổ biến trong lĩnh vực sinh câu hỏi tự động ở giai đoạn trước khi học sâu được áp dụng rộng rãi.

Trong một hướng tiếp cận khác, \textit{Mitkov} và cộng sự \cite{mitkov2006computer} phát triển một hệ thống sinh câu hỏi trắc nghiệm hỗ trợ bằng máy tính, tích hợp nhiều kỹ thuật NLP bao gồm nhận dạng thực thể có tên, gán nhãn từ loại và phân tích cú pháp. Hệ thống này tự động xác định các khái niệm quan trọng trong văn bản nguồn, sau đó áp dụng các quy tắc biến đổi ngôn ngữ linh hoạt để tạo thành câu hỏi và các phương án nhiễu tương ứng.

Các phương pháp truyền thống có ưu điểm nổi bật là khả năng kiểm soát cao về mặt ngữ pháp và cấu trúc câu hỏi đầu ra, đồng thời cho phép người thiết kế can thiệp trực tiếp vào quy trình sinh câu hỏi. Tuy nhiên, các phương pháp này tồn tại nhiều hạn chế căn bản: (i) các câu hỏi được sinh ra chủ yếu nằm ở mức nhận thức thấp theo thang Bloom \cite{bloom1956taxonomy}, tập trung vào khả năng ghi nhớ và tái hiện thông tin bề mặt; (ii) chất lượng của các phương án nhiễu thường không đạt yêu cầu sư phạm do thiếu khả năng suy luận ngữ nghĩa sâu; (iii) hệ thống đòi hỏi thiết kế riêng bộ quy tắc ngôn ngữ học cho từng ngôn ngữ và từng lĩnh vực chuyên môn, gây khó khăn lớn trong việc mở rộng phạm vi ứng dụng. Đặc biệt, các phương pháp này không có khả năng sinh các câu hỏi đòi hỏi suy luận hay phản biện, vốn là các mức nhận thức bậc cao cần thiết trong đánh giá giáo dục hiện đại \cite{anderson2001taxonomy}.

\subsection{Sinh câu hỏi trắc nghiệm dựa trên mô hình ngôn ngữ lớn}

Sự ra đời của các LLM đã tạo ra một bước chuyển quan trọng trong lĩnh vực sinh câu hỏi tự động, từ các hệ thống dựa trên quy tắc sang các phương pháp dựa trên khả năng sinh ngôn ngữ tự nhiên của mô hình. Thay vì thiết kế thủ công các quy tắc biến đổi cú pháp, các nghiên cứu từ năm 2022 trở đi đã tận dụng khả năng hiểu ngữ cảnh sâu và sinh văn bản mạch lạc của LLM để tạo ra câu hỏi trắc nghiệm (MCQ - Multiple-Choice Question) trực tiếp từ đoạn văn bản nguồn.

\textit{Elkins} và cộng sự \cite{elkins2023useful} thực hiện đánh giá có hệ thống về chất lượng câu hỏi trong lĩnh vực giáo dục được sinh bởi các LLM. Kết quả nghiên cứu cho thấy các mô hình này có khả năng tạo ra các câu hỏi đạt yêu cầu sử dụng trong thực tế giảng dạy, tuy nhiên vẫn gặp khó khăn đáng kể trong việc sinh các câu hỏi đòi hỏi tư duy bậc cao theo phân loại Bloom. Nghiên cứu này nhấn mạnh rằng khoảng cách giữa câu hỏi ghi nhớ đơn thuần và câu hỏi phân tích, đánh giá vẫn là thách thức lớn đối với phương pháp sinh câu hỏi chỉ dùng một lần gọi mô hình.

Bên cạnh khả năng sinh câu hỏi, việc sử dụng LLM để tự động đánh giá chất lượng câu hỏi cũng là một hướng đi đáng chú ý. \textit{Moore} và cộng sự \cite{moore2023assessing} đã tiến hành nghiên cứu sử dụng GPT-4 để đánh giá và phát hiện các lỗi biên soạn trong các câu hỏi trắc nghiệm do sinh viên tạo ra. Khi đối chiếu với kết quả đánh giá của chuyên gia giáo dục, nghiên cứu phát hiện ra rằng độ chính xác của GPT-4 (79\%) vẫn kém hơn đáng kể so với một hệ thống đánh giá bằng tập luật truyền thống (91\%). Hơn nữa, GPT-4 có xu hướng ``bắt lỗi'' quá khắt khe, thường xuyên nhận diện sai các lỗi sư phạm không thực sự tồn tại. Phát hiện này củng cố cho nhận định rằng: mặc dù LLM rất mạnh mẽ trong việc xử lý ngôn ngữ, nhưng nếu chỉ sử dụng trực tiếp một LLM đơn lẻ làm giám khảo kiểm định mà thiếu đi các cơ chế kiểm soát chéo, hệ thống sẽ rất dễ đưa ra các đánh giá sai lệch về mặt sư phạm.

\section{Phân tích khoảng trống nghiên cứu}

Dựa trên quá trình tổng quan các công trình liên quan đã trình bày, nghiên cứu rút ra ba khoảng trống chính còn tồn tại trong lĩnh vực sinh câu hỏi trắc nghiệm tự động, thể hiện rõ những hạn chế của các phương pháp hiện hành:

\begin{itemize}
    \item \textbf{Khả năng kiểm soát sư phạm và độ chính xác thấp ở tư duy bậc cao:} Các phương pháp sinh câu hỏi truyền thống dựa trên quy tắc cấu trúc (như phân tích cú pháp NLP, từ điển WordNet) tỏ ra cứng nhắc, chủ yếu tạo ra được những câu hỏi ở mức độ nhận thức thấp. Ngược lại, việc sử dụng trực tiếp một mô hình ngôn ngữ lớn để sinh câu hỏi thường dẫn đến tình trạng ``ảo giác'' (hallucination) hoặc không thể thiết kế các phương án nhiễu có tính chẩn đoán logic. Các mô hình đơn lẻ thường thất bại trong việc phân định ranh giới giữa các cấp độ nhận thức cao như Suy luận và Lập luận phản biện.
    
    \item \textbf{Tốn kém tài nguyên và rào cản triển khai thực tế:} Các hệ thống tiên tiến hiện nay phụ thuộc quá nhiều vào các mô hình thương mại khổng lồ. Dù mang lại chất lượng câu hỏi tương đối tốt, các mô hình này lại đi kèm với chi phí gọi API đắt đỏ, độ trễ hệ thống cao và đòi hỏi kết nối mạng liên tục. Đối với môi trường giáo dục có ngân sách hạn hẹp hoặc các tổ chức yêu cầu bảo mật dữ liệu đề thi nghiêm ngặt, việc phụ thuộc vào nền tảng bên ngoài là một rào cản rất lớn.
    
    \item \textbf{Sự phân mảnh trong quy trình số hóa toàn diện:} Ít có hệ thống nào giải quyết trọn vẹn luồng công việc từ lúc nhận một bức ảnh tài liệu thô (có nhiều yếu tố nhiễu vật lý như mờ, nhòe) đến khi tạo ra bộ câu hỏi hoàn chỉnh. Các công cụ nhận dạng ký tự (OCR) truyền thống thường phá vỡ cấu trúc văn bản, làm mất ngữ cảnh trước khi đưa vào mô hình sinh câu hỏi, khiến luồng dữ liệu trở nên kém hiệu quả.
\end{itemize}

\textbf{Sự cần thiết của giải pháp đề xuất:} Từ những hạn chế cốt lõi trên - sự đắt đỏ và thiếu tính bảo mật của LLM thương mại, sự kém hiệu quả của các quy tắc NLP cũ, cũng như sự thiếu hụt cơ chế tự kiểm duyệt cho các câu hỏi tư duy bậc cao - giải pháp của nghiên cứu này ra đời. Hệ thống sử dụng \textbf{các mô hình cỡ nhỏ mã nguồn mở} nhằm giải quyết triệt để bài toán tài nguyên và bảo mật. Đồng thời, việc ứng dụng kiến trúc \textbf{đa tác tử } đóng vai trò như một cơ chế tự giám sát, mô phỏng quá trình tư duy độc lập, qua đó tối ưu hóa độ chính xác và đảm bảo chất lượng sư phạm cho các câu hỏi trắc nghiệm phức tạp mà không cần sử dụng đến các hệ thống quá khổng lồ.
